\documentclass[conference]{IEEEtran}

\usepackage{cite}
\usepackage{amsmath,amssymb,amsfonts}
\usepackage{algorithmic}
\usepackage{graphicx}
\usepackage{textcomp}
\usepackage{xcolor}
\usepackage{hyperref}

\def\BibTeX{{\rm B\kern-.05em{\sc i\kern-.025em b}\kern-.08em
    T\kern-.1667em\lower.7ex\hbox{E}\kern-.125emX}}

\begin{document}

\title{AI Agent Briefer: A Serverless, Multi-Agent System for Autonomous Daily Intelligence Synthesis\\
}

\author{
\IEEEauthorblockN{Ramadugu Dhanush}
\IEEEauthorblockA{\textit{Dept. of Computer Science and Engineering} \\
\textit{Institute of Aeronautical Engineering (IARE)}\\
Hyderabad, India \\
https://github.com/RDK7159357/ai-agents-using-crewai \\
ramadugudhanush8@gmail.com}
\and
\IEEEauthorblockN{Gautham Sharma}
\IEEEauthorblockA{\textit{Dept. of Computer Science and Engineering} \\
\textit{Institute of Aeronautical Engineering (IARE)}\\
Hyderabad, India \\
gautham.sharma@iare.ac.in}
}

\maketitle

\begin{abstract}
In the modern era, the rapid dissemination of information has led to an overwhelming volume of digital news, making efficient curation a significant challenge. Concurrently, technical interview preparation requires synthesizing vast amounts of decentralized corporate data, which is time-consuming and often yields outdated information. This project presents ``AI Agent Briefer,'' an autonomous, serverless, multi-agent artificial intelligence system designed to deliver high-precision daily technology briefings and in-depth company intelligence. Built on the CrewAI orchestration framework, the system leverages a robust multi-model fallback chain---integrating Large Language Models (LLMs) such as Ollama, Groq, Gemini, OpenRouter, Mistral, and Together AI---to ensure high uptime with zero operational costs. A critical innovation of this project is its rigorous programmatic output validation pipeline, which systematically identifies and mitigates LLM hallucinations, lazy tool execution, and stale data generation. The curated intelligence is delivered via Telegram acting as a multi-modal interface, providing both rich HTML text and synthesized audio via ElevenLabs and Google TTS. This report provides a comprehensive analysis of the system architecture, implementation strategies, validation methodologies, and performance evaluation.
\end{abstract}

\begin{IEEEkeywords}
Multi-Agent Systems, CrewAI, Large Language Models, Serverless Architecture, Output Validation, Generative AI.
\end{IEEEkeywords}

\section{Introduction}

\subsection{Motivation}
The technology industry is characterized by rapid evolution. Professionals and job seekers must continuously monitor advancements in Artificial Intelligence, cybersecurity, startups, and global markets. General-purpose aggregators lack the domain specificity required to filter out noise, while manual research for interview preparation is arduous. Given the advent of autonomous AI agents capable of reasoning, searching the web, and synthesizing data, there is a profound opportunity to automate this knowledge curation.

\subsection{Problem Statement}
Existing AI-driven summarization tools are often hindered by:
\begin{itemize}
    \item \textbf{High API Costs:} Reliance on proprietary models (e.g., GPT-4) prevents continuous daily execution without financial investment.
    \item \textbf{LLM Hallucinations:} Generative models often fabricate dates, summarize outdated articles as ``breaking news,'' or fail to execute web searches, instead writing ``I could not find this information.''
    \item \textbf{Accessibility:} Text-heavy newsletters are difficult to consume in varying environments, highlighting the need for multi-modal (audio) delivery systems.
\end{itemize}

\subsection{Project Objectives}
The primary objectives of this system are:
\begin{enumerate}
    \item To develop distinct AI personas capable of specialized tasks: \textit{News Scout} for daily briefings and \textit{Company Researcher} for interview intelligence.
    \item To engineer a zero-cost inference engine by implementing a dynamic, rate-limit-aware fallback mechanism across multiple free-tier LLM providers.
    \item To design deterministic Python guardrails that programmatically reject subpar generative outputs.
    \item To deploy the system as a serverless infrastructure utilizing GitHub Actions and Telegram webhooks for seamless user interaction via text and audio.
\end{enumerate}

\section{Literature Review}

\subsection{Autonomous Multi-Agent Systems}
Multi-agent systems (MAS) involve multiple interacting intelligent agents that solve problems beyond the capability of a single agent. Frameworks like CrewAI and AutoGen have popularized the use of LLMs as the cognitive engine for these agents, allowing them to collaborate, delegate tasks, and utilize external tools (e.g., search APIs) to gather real-time data autonomously.

\subsection{LLM Hallucinations and Output Validation}
A well-documented limitation of LLMs is ``hallucination''---the generation of fluent but factually incorrect information. In automated news curation, this manifests as stale news being presented as recent. Recent studies suggest that combining stochastic LLM generation with deterministic programmatic validation (e.g., Regex date validation) significantly improves the reliability of enterprise AI systems.

\subsection{Serverless Computing and Microservices}
Serverless computing allows developers to execute code without managing the underlying infrastructure. Platforms like GitHub Actions (for cron jobs) and Vercel (for webhook hosting) enable discrete microservices to operate efficiently, scaling down to zero when idle, which drastically minimizes operational overhead.

\section{System Architecture and Design}

The system is designed with a decoupled architecture, separating orchestration, inference, validation, and delivery pipelines to ensure module independence and high cohesion.

\subsection{Component Architecture}
\begin{itemize}
    \item \textbf{User Interface (Telegram Bot):} Acts as the primary interface. Users can receive automated daily briefs or request on-demand interview prep via \texttt{/brief} or \texttt{/interview <company>}.
    \item \textbf{Orchestrator (\texttt{main.py}):} The central nervous system. It initializes the tasks, sets up the CrewAI environment, and manages the execution flow.
    \item \textbf{Agent Factory (\texttt{agents.py}):} Defines the personas, equips them with tools (Serper Search API), and manages the LLM instantiation logic.
    \item \textbf{Validation Logic:} A series of strict deterministic functions acting as a quality assurance gateway prior to message dispatch.
    \item \textbf{Delivery and TTS Subsystem:} Formats the final validated text into Telegram-friendly HTML and generates an MP3 audio file.
\end{itemize}

\subsection{Data Flow}
1. \textbf{Trigger:} A cron job on GitHub Actions fires daily, or a Telegram webhook is triggered by manual invocation. \\
2. \textbf{Inference:} The Orchestrator requests an LLM from the Agent Factory. The AI Agent receives its prompt and utilizes the Serper tool to fetch live Web Data. \\
3. \textbf{Validation:} The Orchestrator receives the markdown response and parses it through regex-based string validation algorithms. If the output fails (e.g., stale dates, too short, raw tool syntax), the system throws an exception and retries the entire pipeline using the next available LLM in the fallback matrix. \\
4. \textbf{Delivery:} The validated text is sanitized, sent as a Telegram message, and passed to ElevenLabs (or gTTS) for audio synthesis, which is transmitted as an audio file payload.

\section{Implementation Details}

\subsection{The LLM Fallback Chain}
A core innovation of this project is its ability to operate indefinitely on free-tier APIs without failing. The \texttt{get\_llm()} factory function dynamically routes inference requests across different providers. The default priority matrix is:
\begin{enumerate}
    \item \textbf{Ollama (Local)}: Zero network latency and absent API quotas.
    \item \textbf{Groq (Llama 3.3 70B)}: Preferred for extreme inference speed and a generous free-tier allocation.
    \item \textbf{Google Gemini (2.0 Flash)}: Utilized primarily for its massive context window handling intensive search operations.
    \item \textbf{Together AI, OpenRouter, Mistral, HuggingFace}: Kept active as progressive safety-nets.
\end{enumerate}

\textbf{Rate Limit Handling:} The function \texttt{run\_crew\_with\_rate\_limit\_retry()} intercepts HTTP 429 exceptions. It extracts the suggested retry wait time from the payload string, applies exponential backoff, and shifts the state context to gracefully failover to the subsequent provider instead of crashing.

\subsection{Agent Personas and Prompt Engineering}
The system utilizes two highly specialized software agents defined via the CrewAI framework:
\begin{itemize}
    \item \textbf{The News Scout}: Instructed to perform independent searches representing different facets of global and local tech. The system programmatically forces the agent to bypass its innate laziness by stipulating strict demographic coverage (e.g., 3-4 global startups, 3-4 regional Indian stories).
    \item \textbf{The Interview Strategist}: Executes a 4-pronged search methodology (Overview, News, Culture, Challenges). The prompt explicitly commands the generation of ``Killer Questions'' that reference specific facts scraped during the execution process.
\end{itemize}

\subsection{Telegram Message Sanitization}
Because the Telegram API restricts messages to 4096 characters and utilizes an inflexible HTML parser, \texttt{format\_for\_telegram()} acts as a native compiler translating Markdown to Telegram-HTML. It natively splits outputs exceeding character limits across natural paragraph boundaries to maintain structural flow for end-users, while properly substituting Markdown tables and markdown headers to valid list sequences and \verb|<b>| blocks.

\section{Validation and Security Guardrails}

The vulnerability of modern autonomous agents is output reliability. To guarantee production-grade content, this project executes explicit parsing validation protocols on the LLM's raw string payload.

\subsection{News Validation Guardrails}
\begin{itemize}
    \item \textbf{Stale News Regex Filtering:} Extracts geographic date strings (e.g., ``Mar 18 2026'', ``2026-03-18''), converts them into \texttt{datetime} entities, and calculates the delta from the current epoch. If any scraped story exceeds \texttt{NEWS\_MAX\_AGE\_DAYS} (configured to 3), the agent's work is categorically rejected and retried.
    \item \textbf{Deduplication Engine:} Computes an intersectional word overlap between generated headlines. If $\frac{|A \cap B|}{\max(|A|, |B|)} > 0.60$, evaluating two headlines $A$ and $B$, the system registers a deduplication failure, blocking the message dispatch.
    \item \textbf{Topic Diversity Enforcer:} Scans string outputs ensuring non-monochromatic clustering (e.g., preventing a summary composed solely of smartphone peripheral launches) by validating against arrays of mandatory keywords.
\end{itemize}

\subsection{Interview Validation Guardrails}
\begin{itemize}
    \item \textbf{Anti-Lazy Flagging:} If the response string contains the explicit phrase ``Not found in public sources'' repeatedly, the agent is assumed to have omitted internal tool executions. Validation fails, triggering a harsher retrying mechanic.
    \item \textbf{Tool-Call Leakage Prevention:} Prevents raw XML/JSON debug syntaxes (e.g., \verb|<search_query>|, \verb|action_input|) from propagating up the chain into user-facing content.
    \item \textbf{Unfilled Placeholder Detection:} Regex engines sweep the LLM generation for arbitrary brackets indicating unresolved UI templates passed from the instructional prompt directly to output.
\end{itemize}

\section{Results and Performance Evaluation}

\subsection{Cost Efficiency}
By intelligently pivoting request burdens through Groq, OpenRouter, and Gemini Flash's free tiers via \texttt{agents.py}, the operational cost rests permanently at \$0.00/month. Conversely, equivalent autonomous systems navigating web searches on commercial AI architectures like OpenAI's GPT-4o are mathematically estimated to cost upwards of \$40.00/month in daily continuous evaluation.

\subsection{Execution Latency}
\begin{itemize}
    \item \textbf{Daily Briefing:} 25 to 45 seconds total execution time, tightly coupled to network latency during Serper API search executions.
    \item \textbf{Interview Prep:} 40 to 80 seconds. The massive 2-million token context on Gemini handles dozens of raw HTML results seamlessly without chunking architectures, producing exceptional context at the cost of nominal latency spikes.
    \item \textbf{Audio Synthesis:} ElevenLabs API generates 2-minute MP3 deliverables in sub-10s executions, with dynamic automatic fallbacks to generic \texttt{gTTS} to maintain voice inclusion without token expenses.
\end{itemize}

\section{Conclusion and Future Work}

\subsection{Conclusion}
The ``AI Agent Briefer'' successfully demonstrates the viability of deploying highly autonomous knowledge-curation agents on strictly serverless, zero-cost infrastructure models. By acknowledging the innate flaws of modern Large Language Models (specifically hallucinations and functional laziness) and framing them within stringent programmatic control boundaries, the pipeline turns stochastic textual generation into deterministic, enterprise-grade intelligence summaries immediately accessible through mainstream delivery networks such as Telegram.

\subsection{Future Scope}
\begin{enumerate}
    \item \textbf{Vector Database Integration (RAG):} Deploying persistent temporal memory instances via lightweight vector stores (ChromaDB), solving the risk of agents generating the same news event across overlapping days.
    \item \textbf{Dynamic User Profiling:} Deprecating hardcoded geographical logic in favor of a parameterized JSON schema enabling isolated user briefings.
    \item \textbf{Hierarchical Agent Topologies:} Re-orienting the CrewAI pipeline to support parallel execution layers instead of sequential flow, theoretically accelerating execution constraints by an estimated 50\%.
\end{enumerate}

\end{document}
